\documentclass{article}
\usepackage[margin=1in]{geometry}
\usepackage{amsmath,amssymb}
\usepackage[most]{tcolorbox}
\usepackage[hidelinks]{hyperref}

\newcommand{\BookName}{Pattern Recognition and Machine Learning}
\newcommand{\BookAuthor}{Christopher M. Bishop}
\newcommand{\ChapterName}{Introduction}
\newcommand{\ExerciseID}{1.5}

\title{Chapter: \ChapterName}
\author{Ahonakpon Guy Gbaguidi}
\date{}

\begin{document}

\maketitle

\begin{center}
  \large\textbf{Exercise \ExerciseID}

  \vspace{0.5em}

  \normalsize\BookName\ - \BookAuthor
\end{center}

\begin{tcolorbox}[breakable,colback=gray!5,colframe=gray!60,title=Solution]
In this exercise, we have to use the definition (1.38)

\begin{align*}
  \text{Var}[f] = \mathbb{E} \left[ \left( f(x) - \mathbb{E}[f(x)] \right)^2 \right]
\end{align*}

to show that the $\text{Var}[f(x)]$ statifies (1.39):

\begin{align*}
  \text{Var}[f] = \mathbb{E}[f(x)^2] - \left( \mathbb{E}[f(x)] \right)^2
\end{align*}

Let $f(x)$ be a random variable defined on a probability space with probability density function $p(x)$. The expectation of $f(x)$ is given by:

\begin{align*}
  \mathbb{E}[f(x)] &= \sum_x p(x) f(x) \quad \text{(in case of discrete random variable)} \\
  \mathbb{E}[f(x)] &= \int p(x) f(x) dx \quad \text{(in case of continuous random variable)}
\end{align*}

To compute the variance of $f(x)$, we can expand the definition of variance as follows:

\begin{align*}
  \text{Var}[f] &= \mathbb{E} \left[ \left( f(x) - \mathbb{E}[f(x)] \right)^2 \right] \\
  &= \mathbb{E} \left[ f(x)^2 - 2 f(x) \mathbb{E}[f(x)] + \left( \mathbb{E}[f(x)] \right)^2 \right] \\
  &= \mathbb{E}[f(x)^2] - 2 \mathbb{E}[f(x)] \mathbb{E}[f(x)] + \left( \mathbb{E}[f(x)] \right)^2 \\
  &= \mathbb{E}[f(x)^2] - 2 \left( \mathbb{E}[f(x)] \right)^2 + \left( \mathbb{E}[f(x)] \right)^2 \\
  &= \mathbb{E}[f(x)^2] - \left( \mathbb{E}[f(x)] \right)^2
\end{align*}

because the expectation operator is linear, we can separate the terms and compute the expectation 
of each term individually. The final result shows that the variance of $f(x)$ can be expressed as 
the difference between the expected value of $f(x)^2$ and the square of the expected value of 
$f(x)$, which is exactly what we wanted to show.

\par\medskip
\paragraph{Personal note:} I learned more about the properties of the expectation operator, 
such as linearity, and how to manipulate from this MIT OpenCourseWare ressource: 
\url{https://ocw.mit.edu/courses/6-042j-mathematics-for-computer-science-fall-2005/6ad0342f836f80c219470870db432c18_ln14.pdf}.

\end{tcolorbox}

\end{document}
